\lysh{34815}{2}
\begin{alist}
    \item Γνωρίζουμε ότι:
    \[
        \bar{x} = \frac{x_1 \cdot \nu_1 + x_2 \cdot \nu_2 + x_3 \cdot \nu_3 + x_4 \cdot \nu_4}{\nu_1 + \nu_2 + \nu_3 + \nu_4} = \frac{1 \cdot 1 + 2 \cdot 2 + 3 \cdot 1 + 4 \cdot 4}{1+2+1+4} = \frac{1+4+3+16}{8} = \frac{24}{8} = 3
    \]
    \item Για να υπολογίσουμε τη διάμεσο $\delta$, θα πρέπει να βάλουμε τις τιμές σε αύξουσα σειρά: $1, \quad 2, \quad 2, \quad \textbf{3}, \quad \textbf{4}, \quad 4, \quad 4, \quad 4$. Επειδή είναι στο σύνολό τους 8 οι τιμές που έχουμε, η διάμεσος θα ισούται:
    \[
        \delta = \frac{4\text{η παρατήρηση} + 5\text{η παρατήρηση}}{2} = \frac{3+4}{2} = \frac{7}{2} = 3,5
    \]
\end{alist}
\vspace{6mm}
\noindent\rule{\textwidth}{0.4pt}
\vspace{6mm}

